% Generated draft: requires mathematical and source review.
\documentclass[11pt]{article}
\usepackage[margin=25mm]{geometry}
\usepackage{fontspec}
\setmainfont{Latin Modern Roman}
\usepackage{amsmath,amssymb,mathtools,unicode-math}
\setmathfont{Latin Modern Math}
\usepackage{enumitem,needspace,xurl,hyperref}
\setlength{\parindent}{0pt}
\setlength{\parskip}{5pt}
\setlength{\emergencystretch}{3em}
\newcommand{\N}{\mathbb{N}}
\newcommand{\Z}{\mathbb{Z}}
\newcommand{\Q}{\mathbb{Q}}
\newcommand{\R}{\mathbb{R}}
\newcommand{\C}{\mathbb{C}}
\newcommand{\F}{\mathbb{F}}
\newcommand{\A}{\mathbb{A}}
\newcommand{\PP}{\mathbb{P}}
\newcommand{\E}{\mathbb{E}}
\newcommand{\eps}{\varepsilon}
\newcommand{\dd}{\,\mathrm{d}}
\begin{document}
\section*{UnsolvedMath ID: 2745}
\textbf{Source problem code:} catalogue entry
\subsection{Definitions and mathematical statement}
Let \(\Gamma\) be a finite graph. For a tame embedding \(f:\Gamma\hookrightarrow S^3\), use its compact exterior
\[
E(f\Gamma)=S^3\setminus\operatorname{int}N(f\Gamma),
\]
where \(N\) is a regular neighbourhood; this specifies the usual meaning of graph complement in the supplied handlebody discussion. A handlebody is a compact orientable 3-manifold homeomorphic to a regular neighbourhood of a finite connected graph. Following the background, call an embedded graph unknotted when its exterior is a connected sum of handlebodies, allowing the disconnected-subgraph case through that convention. The source does not explicitly fix the embedding category or all empty-graph conventions.

For \(D\subseteq E(\Gamma)\), write \(\Gamma-D\) for deletion of the interiors of the selected edges, retaining their endpoints and hence possible isolated vertices. The intended up-to-\(k\)-deletion reading is
\[
f\text{ is }k\text{-free}\quad\Longleftrightarrow\quad
\forall D\subseteq E(\Gamma),\ |D|\leq k:\
f(\Gamma-D)\text{ is unknotted}.
\]
It includes \(D=\varnothing\). This explicitly interprets the background's duplicated malformed subgraph phrase, without changing the source record.

An unbounded-\(k\) definition is
\[
\operatorname{fr}(\Gamma)=\sup\{k\in\Z_{\geq0}:\exists f:\Gamma\hookrightarrow S^3\text{ that is }k\text{-free}\}.
\]
One must permit \(\infty\) when the admissible \(k\) are unbounded; deletion conditions stabilize once \(k\geq|E(\Gamma)|\). The background calls the invariant an integer but does not state a cap on \(k\) or an infinity convention. Capping at the number of edges would change small-graph readings of the problem, so no such cap is silently imposed.

With a consistent convention for this index, the three supplied assertions are:
\begin{enumerate}
\item Every finite connected cubic graph, meaning degree three at every vertex, has \(\operatorname{fr}(\Gamma)\geq2\). The source does not restrict this part to bridgeless graphs.
\item Every finite graph has \(\operatorname{fr}(\Gamma)\geq2\), including the scope of disconnected graphs intended by the source.
\item There exists a finite graph with \(\operatorname{fr}(\Gamma)=2\), not merely index at least two.
\end{enumerate}
The source defines the invariant first for connected graphs but then asks part (b) for every graph; the extension and small-case conventions need source review.

The supplied background and triage dated 2026-08-17 attribute to Thompson the universal lower bound one, and report \(\operatorname{fr}(K_6)=8\) and Petersen-graph index 4. Neither is an exact-index-two example. They also report index at least two for cubic graphs satisfying the orientable cycle double cover conjecture. For bridgeless cubic graphs the background describes that conjecture as a strong embedding in a closed orientable surface, with complementary regions disks whose closures are disks; it says a sufficiently suitable Heegaard-surface realization in \(S^3\) yields the conditional freeness result. Thus a connected bridgeless cubic graph of index one would refute that conjecture; no such graph is supplied. The background contrasts planar examples and \(K_5\), which can have all subgraphs unknotted, with Conway--Gordon's obstruction for \(K_7\). The registry status is open and the dated triage labels the row partially solved, retaining all three target assertions as unproved in general.
\subsection{Short English statement}
Can every graph be embedded in the 3-sphere so that deleting any two edges still leaves an unknotted spatial graph, and does some graph have exactly this maximum degree of freeness? The cubic case has a supplied conditional result from orientable cycle double covers. The deletion, embedding, disconnected-graph and infinity conventions need explicit source alignment.
\subsection{Sources}
\begin{itemize}
\item Local UnsolvedMath registry, permanent ID 2745, source code catalogue entry, ``Kirby Problem 1.86,'' including supplied literature triage dated 2026-08-17 and the malformed duplicated definition phrase.
\item \emph{K3: A New Problem List in Low-Dimensional Topology}, Problem 1.86 (2026), supplied source: \url{https://math.berkeley.edu/sites/default/files/surv-295-ruberman-watermarked-author-pdf.pdf}.
\item Abigail Thompson, ``The freeness index of a graph,'' arXiv:2206.12939 (2022), source label [Tho22]: \url{https://arxiv.org/abs/2206.12939}.
\item J. H. Conway and C. McA. Gordon, ``Knots and links in spatial graphs,'' \emph{Journal of Graph Theory} 7(4) (1983), 445--453, DOI 10.1002/jgt.3190070410, source label [CG83b].
\item G. Szekeres, ``Polyhedral decompositions of cubic graphs,'' \emph{Bulletin of the Australian Mathematical Society} 8(3) (1973), 367--387, DOI 10.1017/S0004972700042660, source label [Sze73].
\item P. D. Seymour, ``Disjoint paths in graphs,'' \emph{Discrete Mathematics} 29(3) (1980), 293--309, DOI 10.1016/0012-365X(80)90158-2, source label [Sey80].
\end{itemize}
\end{document}
